\documentclass[11pt]{article}
\usepackage[a4paper,margin=26mm]{geometry}
\usepackage[T1]{fontenc}
\usepackage{lmodern}
\usepackage{amsmath,amssymb,amsthm}
\usepackage{microtype}
\usepackage{booktabs}
\usepackage{xurl}
\usepackage{hyperref}
\hypersetup{colorlinks=true,linkcolor=blue,urlcolor=blue,citecolor=blue,
 pdftitle={Connected cographs with distance eigenvalues just below -2},
 pdfauthor={Horace Dong}}
\newtheorem{theorem}{Theorem}
\newtheorem{lemma}{Lemma}
\theoremstyle{remark}
\newtheorem{remark}{Remark}
\newcommand{\repoversion}{1.18.0}
\setlength{\emergencystretch}{1.5em}

\title{Connected cographs with distance eigenvalues\\ just below $-2$}
\author{Horace Dong\\ \small Independent researcher\\
\small\texttt{maxwelldhx@gmail.com}\\
\small ORCID: \href{https://orcid.org/0009-0002-7554-7548}{0009-0002-7554-7548}}
\date{10 November 2026}
\begin{document}
\maketitle

\begin{abstract}
Lou and Lin proved that no connected cograph has a distance eigenvalue in the
open interval $(-2,-1)$. In a recent survey, Manna, Mandal and Tura ask whether
some interval $(-2-\varepsilon,-2)$ is also free of distance eigenvalues of
connected cographs. It is not: for every $m\ge1$ the connected cograph
$K_1\vee(K_1\cup mK_2)$ on $2m+2$ vertices has a distance eigenvalue in
$\bigl(-2-\frac1{2m},-2\bigr)$. The same graphs were used by Conde, Dratman
and Grippo for the adjacency matrix, and the answer also follows in a few lines
from their lemma; we record it because the question is listed as open, with a
direct proof that is formally verified in Lean~4. The proof was found and
formalized by AI agents working under the author's direction; their role is
described at the end of the paper.
\end{abstract}

\section{Introduction}

All graphs are finite and simple. A \emph{cograph} is a graph with no induced
path on four vertices~\cite{cls}. The \emph{distance matrix} of a connected
graph $G$ is $D(G)=(d(u,v))_{u,v}$, where $d$ is the graph distance, and its
eigenvalues are the \emph{distance eigenvalues} of~$G$. Since $D(G)$ is a real
symmetric matrix, they are real.

Lu, Huang and Lou~\cite{lhl} showed that a connected threshold graph has no
distance eigenvalue in $(-2,-1)$, and Lou and Lin~\cite{ll} extended this to
all connected cographs. For threshold graphs, Alazemi, Andeli\'c, Koledin and
Stani\'c~\cite{aaks} found a larger interval below $-1$ that contains no
distance eigenvalue other than~$-2$. Manna, Mandal and Tura~\cite{mmt} observe
that this larger interval does not extend to cographs, with the example
$K_1\vee(K_1\cup3K_2)$, whose distance eigenvalue $\approx-2.155$ lies in it,
and they ask~\cite[Open Problem ``Distance eigenvalues of cographs below $-2$
and sharpness'', part~(i)]{mmt}:
\begin{quote}
Is there $\varepsilon>0$ such that $(-2-\varepsilon,-2)$ contains no distance
eigenvalue of any connected cograph?
\end{quote}

For $m\ge1$ let $G_m=K_1\vee(K_1\cup mK_2)$: a vertex $c$ adjacent to all other
vertices, a vertex $z$ adjacent only to~$c$, and $m$ disjoint edges
$a_ib_i$ ($1\le i\le m$) whose ends are also adjacent to~$c$. The example
of~\cite{mmt} is $G_3$.

\begin{theorem}\label{thm:main}
For every $m\ge1$, the graph $G_m$ is a connected cograph with a distance
eigenvalue in the open interval $\bigl(-2-\frac1{2m},-2\bigr)$. Hence the
answer to the question above is no: for every $\varepsilon>0$ some connected
cograph has a distance eigenvalue in $(-2-\varepsilon,-2)$.
\end{theorem}

The eigenvalue is the smallest root of the cubic $p_m$ of Section~\ref{sec:proof};
numerically it is $-2.383$, $-2.222$, $-2.155$, $-2.049$ and $-2.005$ for
$m=1,2,3,10$ and~$100$.

\paragraph{Relation to the work of Conde, Dratman and Grippo.}
The graphs $G_m$ are the graphs $K_1\vee(mK_2+K_1)$ of Conde, Dratman and
Grippo~\cite[Lemma~6 and its proof]{cdg}, who showed that each has an adjacency eigenvalue
$\lambda_0\in\bigl(0,\frac1{2m+1}\bigr)$, so that positive adjacency eigenvalues
of connected trivially perfect graphs accumulate at~$0$. Their paper does not
discuss the distance matrix, but our answer follows from their lemma in a few
lines (Remark~\ref{rem:cdg}); the survey~\cite{mmt} cites~\cite{cdg}. We record
the answer because the question is listed as open, and we give a direct proof
for $G_m$ itself, which we have formalized.

\begin{remark}\label{rem:cdg}
A connected cograph $G$ on at least two vertices has diameter at most~$2$, so
$D(G)=2J-2I-A(G)$~\cite[Proposition ``Diameter'']{mmt}. Let $x$ be an
eigenvector of $A(G_m)$ for the eigenvalue $\lambda_0$ of~\cite{cdg}, and let
$G=K_1\vee(G_m\cup G_m)$, a connected cograph. The vector $y$ that is $0$ on
the new vertex, $x$ on the first copy and $-x$ on the second satisfies
$A(G)y=\lambda_0y$ and $Jy=0$, so $D(G)y=(-2-\lambda_0)y$, and
$-2-\lambda_0\in\bigl(-2-\frac1{2m+1},-2\bigr)$. This remark is not part of
the formal verification.
\end{remark}

\section{Proof of Theorem~\ref{thm:main}}\label{sec:proof}

\begin{lemma}\label{lem:struct}
$G_m$ is a connected cograph, and for $u\ne v$ we have $d(u,v)=1$ if $u$ and
$v$ are adjacent and $d(u,v)=2$ otherwise.
\end{lemma}

\begin{proof}
Every vertex other than $c$ is adjacent to~$c$, so $G_m$ is connected and any
two non-adjacent vertices are at distance~$2$. Suppose that $v_0v_1v_2v_3$ is an
induced path. If $v_1=c$, then $c$ is adjacent to $v_3$, which is not allowed;
if $v_1=z$, then $v_0=v_2=c$, since $c$ is the only neighbour of~$z$. So
$v_1$ is an end of an edge $a_ib_i$, and its neighbours $v_0,v_2$ lie in
$\{c,a_i,b_i\}\setminus\{v_1\}$. Neither of them is $c$, since $c$ would be
adjacent to the other one; so $v_0=v_2$, a contradiction.
\end{proof}

By Lemma~\ref{lem:struct}, the partition of the vertices into $\{c\}$,
$\{a_1,b_1,\dots,a_m,b_m\}$ and $\{z\}$ is equitable for $D(G_m)$: the row sums
of $D(G_m)$ over the three parts are $(0,2m,1)$ for~$c$, $(1,4m-3,2)$ for each
$a_i$ and~$b_i$, and $(1,4m,0)$ for~$z$. Let
\[
p_m(x)=x^3+3x^2-x-3-m(4x^2+10x+4),
\]
the characteristic polynomial of the quotient matrix
$\left(\begin{smallmatrix}0&2m&1\\1&4m-3&2\\1&4m&0\end{smallmatrix}\right)$.
Then
\[
p_m(-2)=3>0,\qquad
p_m\Bigl(-2-\frac1{2m}\Bigr)=-\frac{4m^2+6m+1}{8m^3}<0,
\]
so by the intermediate value theorem $p_m$ has a root $\lambda$ with
$-2-\frac1{2m}<\lambda<-2$.

We show that $\lambda$ is a distance eigenvalue of $G_m$ with an explicit
eigenvector. Since $\lambda<-2$, $\lambda^2-1\ne0$; put
\[
\alpha=\frac{2m(\lambda+2)}{\lambda^2-1},\qquad \gamma=\lambda\alpha-2m,
\]
and let $x$ be the vector with $x_c=\alpha$, $x_z=\gamma$ and $x_v=1$ for the
other $2m$ vertices; $x\ne0$. By the row sums above,
\begin{align*}
(D x)_c&=2m+\gamma=\lambda\alpha,\\
(D x)_z&=\alpha+4m=\lambda\gamma,\\
(D x)_{a_i}=(D x)_{b_i}&=\alpha+(4m-3)+2\gamma=\lambda.
\end{align*}
The first line is the definition of~$\gamma$. The second is equivalent to
$\alpha(\lambda^2-1)=2m(\lambda+2)$, the definition of~$\alpha$. The third is
equivalent to $\alpha(1+2\lambda)=\lambda+3$, that is, after multiplying by
$\lambda^2-1$, to $2m(\lambda+2)(2\lambda+1)=(\lambda+3)(\lambda^2-1)$, which is
$p_m(\lambda)=0$. Hence $Dx=\lambda x$. Given $\varepsilon>0$, choosing $m$
with $\frac1{2m}<\varepsilon$ proves the last sentence of the theorem.
\qed

\section{Formal verification}\label{sec:lean}

Theorem~\ref{thm:main} is formally verified in Lean~4 (version 4.34.1) with
Mathlib~\cite{mathlib} (version v4.34.1, commit \texttt{d13f23b7}). The files
are in version~\repoversion{} of the repository~\cite{repo}, directory
\texttt{lean/\allowbreak Research/\allowbreak CographDist/}.

\paragraph{The statement.}
The statement file \texttt{Challenge.lean} contains only definitions and the
statement, no proofs. In ASCII notation (implicit arguments omitted):
{\footnotesize
\begin{verbatim}
def HasInducedP4 (G : SimpleGraph V) : Prop :=
  exists f : Fin 4 -> V, Function.Injective f /\
    forall i j : Fin 4, G.Adj (f i) (f j) <-> (i.val + 1 = j.val \/ j.val + 1 = i.val)
def IsCograph (G : SimpleGraph V) : Prop := Not (HasInducedP4 G)
def distMatrix (G : SimpleGraph V) : Matrix V V Real := fun u v => (G.dist u v : Real)
def IsDistEigenvalue [Fintype V] (G : SimpleGraph V) (mu : Real) : Prop :=
  exists x : V -> Real, x != 0 /\ (distMatrix G).mulVec x = mu * x
def OpenProblemI_negative : Prop :=
  forall eps : Real, 0 < eps -> exists (V : Type) (_ : Fintype V) (G : SimpleGraph V),
    G.Connected /\ IsCograph G /\
    exists mu : Real, -2 - eps < mu /\ mu < -2 /\ IsDistEigenvalue G mu
\end{verbatim}}
\noindent (In the file, \texttt{mu * x} is the scalar multiple
\texttt{mu • x}, and \texttt{distMatrix} is marked \texttt{noncomputable}.) Mathlib's \texttt{SimpleGraph.dist} is $0$ for vertices in
different components; the statement asks for connected graphs. A test file
checks that the path on four vertices is not a cograph, that $K_2$ is a
connected cograph, and that $1$ is and $2$ is not a distance eigenvalue
of~$K_2$.

\paragraph{Review and freezing.}
An independent Claude Code agent, which had not written the file, reviewed the
statement and the tests against the definitions and the question
of~\cite{mmt} and approved them; its three minor suggestions (a comment and
three more tests) were applied, and the file was frozen on 9 November 2026 at
22:56~UTC by recording the SHA-256 hash
\begin{center}\small
\texttt{960a24b9de500d2449a73a83b7385653f63054ab52dab3053d287c4e1270e569}
\end{center}
of \texttt{Challenge.lean}. The proof files were drafted while the review was
running; the statement file was not changed to fit them, and after the review
only its comment changed. Moving the files to the directory of the repository
changed only their \texttt{import} lines; \texttt{Challenge.lean} itself,
which imports only Mathlib, did not change.

\paragraph{The proof.}
Three files prove the theorem along the lines of Section~\ref{sec:proof}:
\texttt{Family.lean} (the graph $G_m$, its connectedness, its distances and the
absence of an induced $P_4$), \texttt{Eigen.lean} (the root of $p_m$ by the
intermediate value theorem, and the eigenvector) and \texttt{Main.lean}
(\texttt{openProblemI\_negative}, and \texttt{family\_theorem}, the statement of
Theorem~\ref{thm:main} for each~$G_m$). The final audit
\texttt{Audit.lean} restates both and prints their axioms and those of the
five tests: each is exactly \texttt{propext}, \texttt{Classical.choice} and
\texttt{Quot.sound}. No file uses \texttt{sorry}, \texttt{native\_decide},
\texttt{bv\_decide} or an \texttt{axiom} declaration.

\paragraph{Checks.}
On 9 November 2026 (23:20--23:23~UTC) the six files were compiled, one module at
a time, in two new empty output directories, one after the other, on the
author's Mac, with the pinned toolchain and Mathlib, which were not rebuilt;
each module was then replayed with \texttt{leanchecker}. All $24$ commands
ended with exit code~$0$, with no error, no warning and no \texttt{sorry}, and
the compiled \texttt{.olean} files of the two directories are byte for byte
identical. From 9 to 10 November 2026 the whole Lean project of the
repository, with these six files, was rebuilt from a new empty directory on the
author's Windows computer, one module at a time and not with \texttt{lake}: all
$3005$ modules compiled and $3004$ passed \texttt{leanchecker}; the remaining
one, the root of the project, has only imports and no constants. The formalization has not been checked with comparator or with
an independent kernel.

\section*{How this paper was produced}
This work was carried out on 9 November 2026 (UTC) in a workflow directed by the
author, using agents of Anthropic Claude Code with the model Claude Opus~5.5.
A Claude Code agent scouting for open problems found the question, noticed that
the family $G_m$ answers it, and checked this numerically and symbolically.
Another Claude Code agent searched the literature (22:49--23:17~UTC) and found
the lemma of~\cite{cdg} and Remark~\ref{rem:cdg}. Claude Code wrote the
statement file, the proofs and this text and ran the checks of
Section~\ref{sec:lean}. An independent Claude Code agent (model Claude
Fable~5.1), which had written none of the files, reviewed the statement file
and compiled it with its tests; before the release a second such agent (same
model) reviewed this paper, the Lean files and the records of the builds,
without running Lean. The corrections of both reviews were applied. OpenAI Codex was not involved.
Throughout, ``we'' refers to this workflow; the reading and searching described
in the paper were done by the agents. The author is not a professional
mathematician.
The author has read this paper and, using a side-by-side table, compared its Lean
statement with Theorem~\ref{thm:main}.
No human mathematician has reviewed this paper. The correctness of the result
rests on the Lean proof and on the agreement of the Lean statement with the
question of~\cite{mmt}, which was checked by the reviewing agent; the written
proof is given for the reader. AI systems are not authors of this paper; the
author takes full responsibility for its content.

\medskip
\noindent Thanks are due to Pallabi Manna, Santanu Mandal and Fernando C.\ Tura
for the question, and to C.~M.~Conde, E.~Dratman and L.~N.~Grippo, whose lemma contains the
essential example.

{\small
\begin{thebibliography}{99}
\bibitem{aaks}
A.~Alazemi, M.~Andeli\'c, T.~Koledin and Z.~Stani\'c, Eigenvalue-free
intervals of distance matrices of threshold and chain graphs, \emph{Linear
Multilinear Algebra} \textbf{69} (2021), 2959--2975,
\url{https://doi.org/10.1080/03081087.2019.1701624}.
\bibitem{cdg}
C.~M.~Conde, E.~Dratman and L.~N.~Grippo, A sharp forbidden interval for the
nontrivial adjacency eigenvalues of trivially perfect graphs,
arXiv:2606.06052v1, 4 June 2026.
\bibitem{cls}
D.~G.~Corneil, H.~Lerchs and L.~Stewart Burlingham, Complement reducible
graphs, \emph{Discrete Appl.\ Math.} \textbf{3} (1981), 163--174.
\bibitem{repo}
H.~Dong, \emph{ai4math-results}, repository,
\url{https://github.com/Horace-Maxwell/ai4math-results}, version~\repoversion{}
(2026); archived at Zenodo under the concept DOI
\url{https://doi.org/10.5281/zenodo.22970254}.
\bibitem{ll}
Z.~Lou and H.~Lin, Distance eigenvalues of a cograph and their multiplicities,
\emph{Linear Algebra Appl.} \textbf{608} (2021), 1--12,
\url{https://doi.org/10.1016/j.laa.2020.08.016}.
\bibitem{lhl}
L.~Lu, Q.~Huang and Z.~Lou, On the distance spectra of threshold graphs,
\emph{Linear Algebra Appl.} \textbf{553} (2018), 223--237,
\url{https://doi.org/10.1016/j.laa.2018.05.014}.
\bibitem{mmt}
P.~Manna, S.~Mandal and F.~C.~Tura, The spectral theory of cographs: a survey,
organized by graph matrix, arXiv:2611.06981v1, 4 November 2026; we quote its
open problem and use its names of results.
\bibitem{mathlib}
The mathlib Community, The Lean mathematical library, in \emph{Proc.\ 9th ACM
SIGPLAN Int.\ Conf.\ Certified Programs and Proofs} (CPP '20), ACM, 2020,
\url{https://doi.org/10.1145/3372885.3373824}.
\end{thebibliography}}
\end{document}
